\documentclass[11pt]{article}
\usepackage[T5,T1]{fontenc}
\usepackage[utf8]{inputenc}
\usepackage[vietnamese,english]{babel}
\usepackage[margin=1in]{geometry}
\usepackage{amsmath,amssymb}
\usepackage{enumitem}
\usepackage{microtype}

\ifdefined\pdfcatalog
  \pdfcatalog{/Lang (en-US)}
\fi

\setlength{\parindent}{0pt}
\setlength{\parskip}{0.55em}
\setlength{\emergencystretch}{2em}
\setlist{nosep,leftmargin=*}

\newcommand{\findinglabel}[1]{\textbf{#1}}
\newcommand{\classification}[1]{\textit{Classification: #1.}}

\title{Technical Review of ``Early-Stopping Activation Steering for Hallucination Mitigation in Vietnamese Clinical Language Models''}
\author{Manuscript review}
\date{August 23, 2026}

\begin{document}
\selectlanguage{english}
\maketitle

\section*{Overall assessment}

The manuscript has addressed several earlier reproducibility and interpretation problems.  It now distinguishes BERTScore reference preference from clinician-adjudicated correctness, defines the $D_1$-matched control, discloses the $N=499$ analysis population, reports paired tests for the main proxy comparison, and qualifies the cumulative-dose confound.  These are substantive improvements.  Important problems nevertheless remain: the title and abstract still frame a semantic proxy as hallucination suppression; the natural-termination experiment reduces the effect to an untested $+1.00$ percentage point; the ``equal-initial-$\alpha$'' matrix includes a matched-dose row that does not share the initial $\alpha_0$; the ablation and layer-sweep differences have no uncertainty analysis; the dataset source is cited with the wrong edition; and the source does not compile.  The work is promising but still requires major revision before clinical or causal claims are justified.

\section*{Major technical findings}

\subsection*{1. The title-level hallucination claim still exceeds the proxy evidence}

\classification{Unsupported central claim}

\findinglabel{Location.} Title; Abstract; Introduction; contribution~3; BERTScore criterion; Conclusion.

\findinglabel{Problem.} The manuscript correctly states that reference-preference accuracy is a semantic proxy and not clinically adjudicated correctness.  However, the title still claims ``Hallucination Mitigation,'' and the Abstract says the framework ``dynamically suppress[es] hallucinations'' and prevents representation saturation.  The evaluation only establishes whether an output is closer in BERTScore to one reference $y^+$ than to one constructed counter-answer $y^-$.  An output may win that comparison while containing a wrong dose, omitted contraindication, or a different unsafe claim.

\findinglabel{Why it matters.} A limitation acknowledged later does not make the stronger title and abstract claims valid.  Clinical hallucination mitigation requires direct factual/harm labels, not only similarity to two templates.

\findinglabel{Suggested fix.} Either add blinded clinician adjudication with factual-correctness and harm endpoints, or narrow the title and headline claims to ``BERTScore Reference-Preference and Repetition.''  Report the proxy as preliminary evidence only.  Validate its sensitivity, specificity, and failure modes against clinician labels, especially for numbers, units, negation, pregnancy status, and drug interactions.

\subsection*{2. The natural-termination experiment does not confirm a persistent binary benefit}

\classification{Definite statistical overstatement}

\findinglabel{Location.} ``Natural Termination Validation'' in Discussion; Abstract's main $+4.21$-point result.

\findinglabel{Problem.} At 800 tokens, reference-preference accuracy changes from $64.80\%$ to $65.80\%$ over 500 questions, corresponding to 324 versus 329 reference-preferred outputs: only five net additional successes.  No McNemar table, confidence interval, or $p$-value is reported.  For a two-sided exact McNemar test, even the most favorable possible discordant table---zero regressions and five improvements---has $p=0.0625$; any additional discordant pairs make the evidence weaker.  The accompanying BERTScore change from $0.6804$ to $0.6834$ also has no paired interval or test.

\findinglabel{Why it matters.} The sentence ``confirming that steering benefits persist'' is not supported for the binary proxy.  More importantly, the effect shrinks from $+4.21$ points under forced 200-token generation to $+1.00$ point under natural termination, indicating that truncation materially affects the headline result.

\findinglabel{Suggested fix.} Move the 800-token experiment into Methods and Results, provide the full paired transition table, exact McNemar result, BERTScore paired effect and CI, output-length distribution, and category results.  Use naturally terminated generation as the primary clinically relevant analysis.  Describe the 200-token result as a fixed-length stress test and replace ``confirming'' with descriptive wording unless the natural-termination effect is statistically and clinically supported.

\subsection*{3. The main paired statistics are not fully reproducible}

\classification{Implementation and reporting ambiguity}

\findinglabel{Location.} Abstract; BERTScore Reference-Preference Results; ``Effect Sizes and Uncertainty.''

\findinglabel{Problem.} The $N=499$ marginal rates imply 365 baseline and 386 steered reference-preferred outputs, but the manuscript does not report the discordant-pair counts required to reproduce McNemar's $p=0.006$, nor whether an exact, continuity-corrected, or asymptotic test was used.  Ties are defined as neither class, but their counts and encoding in McNemar's test are omitted.  The Wilcoxon signed-rank result ($p=0.017$) is paired with a ``95\% CI for mean difference,'' although Wilcoxon targets a rank/location contrast; the CI construction and the actual baseline/steered mean BERTScore values are not reported.

\findinglabel{Why it matters.} A plausible $p$-value is not enough for reproduction.  Test choice, tie handling, effect estimate, and interval method determine the interpretation.

\findinglabel{Suggested fix.} Report the $2\times2$ paired table, tie counts, test variant, alternative hypothesis, and software call.  Give the point estimate of the paired BERTScore difference and state whether its CI is paired-bootstrap, $t$-based, or another method.  If Wilcoxon is retained, report the rank-based effect or Hodges--Lehmann estimate separately from a bootstrap CI for the mean.  Predefine primary outcomes and multiplicity handling.

\subsection*{4. The ``equal-initial-$\alpha$'' matrix does not have equal initial magnitude in all four rows}

\classification{Definite terminology and design-description error}

\findinglabel{Location.} Contribution~4; Methodology's matched-dose definition; Equal-Initial-$\alpha$ Results and table.

\findinglabel{Problem.} Continuous, hard-cutoff, and linear-decay schedules begin at $\alpha_0$.  The matched-dose control instead applies
\[
\alpha_{\mathrm{match}}=D_{\mathrm{decay}}/T
=\frac{K+1}{2T}\alpha_0.
\]
With $K=16$ and $T=200$, its actual constant coefficients are $0.6375$, $0.765$, and $0.85$ for source values $\alpha_0=15$, 18, and 20.  Therefore the matched-dose row does not share the stated initial magnitude and should not be described as the fourth schedule in a matrix ``holding initial steering strength constant.''  The table's first column labels that row with $\alpha_0$, although the applied coefficient is $\alpha_{\mathrm{match}}$.

\findinglabel{Why it matters.} The experiment contains two distinct control families: three equal-initial-scale schedules and a separate equal-$D_1$ control.  Conflating them obscures which variable is held constant.

\findinglabel{Suggested fix.} Present a $3\times3$ equal-initial-$\alpha_0$ comparison plus a separately labeled $D_1$-matched control.  Add an ``actual per-token coefficient'' column showing $\alpha_{\mathrm{match}}$.  Continue to state that $D_2=\sum_t\alpha(t)^2$ and temporal distribution remain unmatched.  Use ``cumulative magnitude'' rather than ``energy'' for $D_1$.

\subsection*{5. Fine-grained schedule rankings lack uncertainty}

\classification{Unsupported comparative claim}

\findinglabel{Location.} Abstract; contribution~3; Equal-Initial-$\alpha$ Results and table; Conclusion.

\findinglabel{Problem.} The paper reports no confidence intervals or paired tests for the ablation table.  The differences supporting the schedule claim are small: at $\alpha_0=18$, hard cutoff and linear decay exceed continuous steering by $0.0018$ and $0.0017$ BERTScore F1, while hard cutoff lowers Rep-4 by only $0.01$ percentage point.  Twelve schedule/intensity conditions and several metrics are inspected, but no prespecified contrast or multiplicity control is given.

\findinglabel{Why it matters.} Raw rank ordering does not establish that early stopping reliably outperforms continuous steering.  The best schedule also changes by metric and $\alpha_0$, and ROUGE-L consistently favors continuous steering.

\findinglabel{Suggested fix.} Report paired differences and 95\% CIs for prespecified early-stop versus continuous comparisons at each $\alpha_0$, with multiplicity control or an explicitly exploratory designation.  Avoid ``achieve higher'' as a general finding if intervals include zero.  State the multi-objective trade-off rather than selecting a winner from tiny untested differences.

\subsection*{6. Layer-selection evidence is incomplete and partly misstated}

\classification{Definite arithmetic error and unsupported representation claim}

\findinglabel{Location.} Contribution~2; ``Layer-Wise Subspace Probing'' in Discussion; Methodology.

\findinglabel{Problem.} Only a range, $0.7040$--$0.7140$, is reported for layers 4--15; the per-layer results, Layer~8 value, uncertainty, and selection criterion are absent.  The displayed range has width $0.0100$, so the statement $\Delta<0.01$ is false as written.  The text calls the experiment ``subspace probing,'' but describes sweeping activation steering across layers, which measures downstream intervention performance rather than linear probe separability.  It then infers that truthfulness representations are distributed across layers, a representation-level conclusion that does not follow from the reported intervention-score range alone.

\findinglabel{Why it matters.} Readers cannot verify why Layer~8 was selected or distinguish a probe result from a steering ablation.  The same $N_{eval}=100$ subset may also have been used for tuning, so ``held-out'' is ambiguous.

\findinglabel{Suggested fix.} Add a per-layer validation table with the exact metric, uncertainty, and selected tie rule.  Call it a steering-layer sweep unless a separate trained probe is reported.  Identify the subset as validation or test, not merely held-out, and do not infer distributed truthfulness representations without direct representation analysis.

\subsection*{7. Vector extraction may encode completed-answer artifacts rather than an early-generation truthfulness direction}

\classification{Likely representation confound}

\findinglabel{Location.} ``Contrastive Representation Extraction'' and ``Early-Stopping Steering Hook.''

\findinglabel{Problem.} The vector is estimated at the final token of the complete concatenated sequence $x_i\oplus y_i^{\pm}$, after the model has seen the entire truthful or hallucinated answer.  It is then injected into the first 16 generated-token states, which have a different context and temporal role.  Synthetic counter-answers may also contain lexical or template cues at their final token.  No test shows that the resulting direction generalizes from answer-final states to early-generation states or to independently authored errors.

\findinglabel{Why it matters.} The intervention can work empirically without the direction being a general truthfulness representation.  The current mechanism claim may instead reflect response-style or perturbation-template artifacts.

\findinglabel{Suggested fix.} Compare final-token extraction with prompt-only, first-answer-token, and pooled answer representations.  Keep all variants of each source/question in one split, test template- and source-disjoint transfer, and evaluate independently written hallucinations.  Report a control direction based on matched lexical/style contrasts.

\subsection*{8. Dataset provenance and splitting remain insufficient, and the cited edition is wrong}

\classification{Definite bibliographic error and likely leakage risk}

\findinglabel{Location.} Abstract; contribution~1; Experimental Setup; bibliography entry \texttt{ref6}.

\findinglabel{Problem.} The manuscript calls the source the Vietnamese National Drug Formulary, 5th Edition, 2018.  The 2018 formulary is identified as the \emph{second} edition; a third edition was subsequently issued in 2022.  The fifth-edition label appears to be inherited from the distinct Vietnamese Pharmacopoeia and is not valid for the Drug Formulary.  The paper also gives no dataset-construction section describing source passages, Q\&A generation, hallucinated counter-answer construction, quality control, deduplication, licensing, or release.  A 70/15/15 split of 14,700 yields 2,205 nominal test records, yet only 500 are sampled without a sampling rule or seed.  ``Question-disjoint'' does not guarantee monograph-, drug-, or template-disjoint evaluation.

\findinglabel{Why it matters.} Incorrect source identity undermines auditability, and source/template overlap can inflate vector transfer and BERTScore preference on the test set.

\findinglabel{Suggested fix.} Correct \texttt{ref6} to the exact edition actually used and cite the official issuance record.  Add a dataset card with provenance, item-generation rules, expert review if any, licenses, category counts, split counts, sampling seed, and source identifiers.  Group splits by source monograph/drug and generation template before producing positive/negative variants, and quantify cross-split similarity.

\subsection*{9. Fixed-length and natural-generation protocols are incompletely connected}

\classification{Evaluation-design ambiguity}

\findinglabel{Location.} Experimental Setup; ablation table; Natural Termination Validation; Abstract.

\findinglabel{Problem.} The 200-token evaluation deliberately truncates almost every response, while the 800-token experiment allows natural completion.  The manuscript does not report mean/median lengths, finish-reason distributions, prompt lengths, or whether reference answers have comparable lengths.  The main abstract effect comes from the fixed-length condition even though the natural condition produces a substantially smaller result.  The 800-token experiment is described only in Discussion and lacks a table, decoding details, repetition, ROUGE-L, latency, and category breakdown.

\findinglabel{Why it matters.} Long irrelevant continuations and truncation can change every reported similarity and repetition metric.  A clinical deployment claim should be based on realistic complete answers, not a fixed-length stress condition.

\findinglabel{Suggested fix.} Define one primary generation protocol with natural stopping and a clinically appropriate maximum.  Report length and finish-reason distributions and all primary metrics under that protocol.  Retain the 200-token experiment only as a secondary stress test and explain why fixed-length output is scientifically useful.

\subsection*{10. Metric and hook implementations remain incomplete}

\classification{Implementation ambiguity}

\findinglabel{Location.} Methodology and Experimental Setup.

\findinglabel{Problem.} Rep-4 is not mathematically defined (macro versus micro averaging, repeated $n$-gram denominator, tokenizer, and treatment of short outputs are unknown).  ROUGE-L tokenization for Vietnamese is unspecified.  BERTScore identifies model and layer but omits IDF use, baseline rescaling, package/version, and numerical tie tolerance.  The hook description omits the exact Qwen module, zero/one-based layer indexing, prompt-forward handling, whether the vector is added to all positions or only $[:, -1, :]$, KV-cache behavior, dtype/device placement, and hook removal.  Hardware count, batch size, software versions, random seeds, and latency methodology are absent.

\findinglabel{Why it matters.} These choices can materially change the metrics, intervention timing, and runtime.  ``No consistent latency penalty'' is not statistically evaluable from one mean per condition.

\findinglabel{Suggested fix.} Provide explicit metric formulas and complete library configurations, release evaluation code, and include minimal hook pseudocode.  Report synchronized repeated latency measurements after warm-up with batch size, prompt/output lengths, mean, and dispersion.  Replace the broad ``architecture-agnostic'' claim with ``conceptually applicable to other residual-stream transformers, subject to architecture-specific validation.''

\section*{Equation, notation, sign, branch, and dimensional audit}

The mean-difference direction $\mu_l^+-\mu_l^-$, unit normalization, and addition of $+\alpha(t)v_{steer}^{(l)}$ are sign-consistent, and the hidden-state/vector dimensions are compatible.  The formulas for $D_{\mathrm{continuous}}$, $D_{\mathrm{cutoff}}$, and $D_{\mathrm{decay}}=(K+1)\alpha_0/2$ are arithmetically correct for nonnegative $\alpha_0$ and the stated endpoint $\alpha(K)=\alpha_0/K$.  No inverse-trigonometric, coordinate-transform, or branch/quadrant expressions occur.  The remaining implementation-facing notation issue is role drift: $\alpha_0$ means the actual initial coefficient for three schedules but only indexes the source decay dose for the matched row, whose applied coefficient is $\alpha_{\mathrm{match}}$.  The concatenation operator $\oplus$, token position, chat template, and generation-step counter also require explicit definitions.

\section*{Meaningful editorial, compilation, and reference findings}

\subsection*{The source still does not compile with pdfLaTeX}

\findinglabel{Location.} Bibliography entry \texttt{ref6}.

\findinglabel{Problem.} A clean \texttt{pdflatex} run stops at \texttt{ref6} because commands such as \verb|\uhorn| are undefined under the current OT1-only setup.  The source reaches the bibliography but produces no PDF from a clean build.

\findinglabel{Why it matters.} The manuscript is not submission-ready, and the Vietnamese title cannot be verified typographically in the generated document.

\findinglabel{Suggested fix.} Use a venue-compatible Unicode engine or configure pdfLaTeX with T5/Vietnamese encoding, then enter the verified Vietnamese title consistently.  Rebuild twice from a clean directory and inspect all citation, table-width, and font warnings.

\subsection*{Citation fit remains uneven}

\findinglabel{Location.} Introduction and bibliography.

\findinglabel{Problem.} \texttt{ref10} is a LoRA paper rather than a direct source for the broad SFT mitigation claim; \texttt{ref12} establishes catastrophic forgetting generally rather than loss of LLM instruction-following ability after the described SFT; and \texttt{ref11} does not directly establish that ``internal activation pathways'' cause RAG models to ignore evidence.  The dataset source entry is factually mis-editioned as noted above.

\findinglabel{Why it matters.} Indirect citations blur the boundary between prior evidence and the manuscript's hypotheses.

\findinglabel{Suggested fix.} Cite studies directly evaluating SFT/DPO for hallucination, LLM forgetting after domain adaptation, and evidence-use failures in RAG.  Rewrite mechanistic language as a hypothesis when the cited work does not test that mechanism.

\section*{Proofing sweep}

The bundled mechanical scan was run on both the current source and the available PDF and returned no high-confidence pattern hits.  The bibliography was checked manually for duplicate punctuation and malformed quotation marks; no such micro-defects were found.

The following bounded manual corrections are high-confidence:
\begin{itemize}
    \item \textbf{Abstract/contribution~5:} the displayed $73.15\%$ and $77.35\%$ subtract to $4.20$ points, while the count-derived unrounded difference is $21/499\approx4.21$ points.  Report the counts or state that deltas use unrounded rates.
    \item \textbf{Category table:} the same rounding issue produces displayed differences of $5.96$ and $2.88$ for Pregnancy and Dosage, versus count-derived $5.95$ and $2.87$.  Add counts or calculate display-consistent deltas.
    \item \textbf{Layer sweep:} change $\Delta<0.01$ to $\Delta=0.0100$ based on the displayed endpoints, or report sufficient precision to justify a strict inequality.
    \item \textbf{Matched Dose row:} display $\alpha_{\mathrm{match}}$ rather than implying that the applied initial coefficient equals the row's $\alpha_0$.
    \item \textbf{Terminology:} replace ``zero-training'' with ``no model-weight updates'' or ``no model retraining,'' since vector estimation, validation, and layer selection require labeled data and computation.
\end{itemize}

\section*{Items requiring external verification}

\begin{itemize}
    \item Verify the exact National Drug Formulary edition, official Vietnamese/English title, issuance record, year, publisher, and reuse permissions; the stated fifth edition is inconsistent with the documented 2018 second edition.
    \item Verify every dosage, contraindication, pregnancy/lactation, and interaction reference and counter-answer against current authoritative clinical sources and clinician review.
    \item Validate the BERTScore proxy and the selected multilingual BERT layer for Vietnamese clinical negation, numbers, units, and entailment before using it as evidence of factual alignment.
    \item Verify novelty against current token-dependent and schedule-based activation-steering work and test the claimed applicability beyond Qwen2.5-7B-Instruct.
    \item Verify privacy and resource-constrained deployment claims with an explicit threat model, memory footprint, throughput, and hospital-infrastructure assumptions.
\end{itemize}

\section*{Highest-priority revision sequence}

\begin{enumerate}
    \item Make naturally terminated generation the primary evaluation and correctly test/report the 800-token results.
    \item Obtain clinician-adjudicated correctness and harm labels or narrow the title and abstract to the BERTScore proxy.
    \item Report complete paired statistics, uncertainty, and prespecified schedule contrasts, including ablation CIs.
    \item Separate the equal-initial-$\alpha_0$ experiment from the equal-$D_1$ matched control and display the actual $\alpha_{\mathrm{match}}$ values.
    \item Add per-layer results, clarify the validation split, and test whether final-answer vectors transfer to early-generation states without template leakage.
    \item Correct the National Drug Formulary edition and add complete dataset provenance, source-grouped splits, and counter-answer construction details.
    \item Complete metric/hook specifications, correct numerical/rounding statements, repair the Vietnamese LaTeX encoding, and strengthen citation fit.
\end{enumerate}

\end{document}